\documentclass[11pt,a4paper]{article}
\usepackage[utf8]{inputenc}
\usepackage[T1]{fontenc}
\usepackage[spanish,es-tabla]{babel}
\usepackage{lmodern}
\usepackage[a4paper,margin=2.5cm,headheight=15pt]{geometry}
\usepackage{microtype}
\usepackage{amsmath,amssymb}
\usepackage{booktabs,tabularx,longtable,array}
\usepackage{graphicx,xcolor}
\usepackage{pgfplots}
\pgfplotsset{compat=1.18}
\usepackage{tikz}
\usetikzlibrary{arrows.meta,positioning}
\usepackage{caption}
\usepackage{placeins}
\usepackage{enumitem}
\usepackage[round,authoryear]{natbib}
\usepackage{fancyhdr}
\usepackage{lastpage}
\usepackage{hyperref}
\hypersetup{colorlinks=true,linkcolor=black,citecolor=black,urlcolor=teal,
 pdftitle={Violencia familiar y seguridad ciudadana en El Salvador: MICS e Infosegura},
 pdfauthor={Borrador para discusión académica}}
\definecolor{petroleo}{HTML}{146475}
\definecolor{gris}{HTML}{52606B}
\definecolor{ocre}{HTML}{B47530}
\newcolumntype{Y}{>{\raggedright\arraybackslash}X}
\newcommand{\var}[1]{\texttt{\detokenize{#1}}}
\newcommand{\fuente}[1]{\par\vspace{3pt}{\footnotesize\textit{Fuente y alcance:} #1\par}}
\newcommand{\pendiente}{\textit{Pendiente}}
\setlength{\parindent}{1.2em}
\setlength{\parskip}{3pt}
\renewcommand{\arraystretch}{1.18}
\captionsetup{font=small,labelfont=bf,skip=7pt}
\setlist{nosep,leftmargin=*}
\pagestyle{fancy}
\fancyhf{}
\fancyhead[L]{\small MICS 2014 e Infosegura 2021}
\fancyhead[R]{\small Violencia y protección}
\fancyfoot[C]{\small \thepage\ de \pageref{LastPage}}
\renewcommand{\headrulewidth}{0.3pt}
\emergencystretch=2em
\makeatletter
\setlength{\@fptop}{0pt}
\setlength{\@fpsep}{16pt plus 2pt minus 2pt}
\setlength{\@fpbot}{0pt plus 1fil}
\makeatother
\title{\vspace{-1.8cm}\includegraphics[width=0.72\textwidth]{images/logo_infosegura_usaid_pnud.png}\\[-0.2cm]
\textbf{Violencia familiar y seguridad ciudadana en El Salvador}\\[0.4em]
\large Disciplina infantil, normas de género y uniones tempranas\\
en diálogo con el análisis multidimensional de Infosegura}
\author{Borrador para discusión académica}
\date{Noviembre de 2021}

\begin{document}
\maketitle
\thispagestyle{fancy}
\begin{abstract}
La seguridad ciudadana requiere observar tanto los delitos registrados como
las relaciones familiares y las condiciones de vida. Este artículo articula
los módulos de disciplina infantil, actitudes de género y primera unión de la
ENS MICS 2014 con la presentación de Infosegura sobre El Salvador de 2021.
Combina una auditoría del cuestionario, tabulados oficiales y agregados de
homicidios, victimización y contexto económico. MICS publica 52,0\,\% de
disciplina violenta en niños de 1 a 14 años. Infosegura presenta una reducción
de la tasa total de homicidios de 106,3 en 2015 a 20,9 en 2020 por cada cien mil
habitantes, mientras su proyección de 2021 muestra señales diferentes para
las muertes de mujeres. La lectura conjunta amplía el diagnóstico, sin
identificar una trayectoria individual entre disciplina y violencia letal.
Se preservan las unidades, períodos y límites de cada fuente, incluidas las
proyecciones y las bases incompletamente documentadas. Las cifras son
secundarias, no reestimaciones propias. El protocolo en R considera diseño
complejo y desconocidos y se verificó con datos artificiales; la aplicación a
microdatos reales permanece pendiente. Se propone una agenda que vincula
apoyo a cuidadores, prevención de violencia de género y oportunidades de desarrollo.
\end{abstract}
\noindent\textbf{Palabras clave:} violencia familiar; disciplina infantil;
normas de género; unión temprana; seguridad ciudadana; Infosegura; El Salvador.

\section{Introducción}
La violencia en el hogar es un problema de derechos y de medición. Los hechos
conocidos por instituciones no coinciden necesariamente con las experiencias
que una encuesta puede captar. El registro administrativo depende de la
denuncia, el acceso a servicios y las prácticas de documentación; una encuesta
depende de su población cubierta, sus preguntas, su informante y la disposición
a responder. Ninguna fuente representa por sí sola el universo del fenómeno.

En este contexto, la Encuesta Nacional de Salud de Indicadores Múltiples por
Conglomerados de 2014, en adelante ENS MICS, permite estudiar tres dimensiones
relacionadas: las prácticas disciplinarias dirigidas a un niño seleccionado,
las opiniones de las mujeres frente a situaciones de violencia contra la esposa
y la edad de inicio de la primera unión. Su valor analítico depende de mantener
claras sus diferencias conceptuales y temporales.

El objetivo es construir una estrategia descriptiva verificable. Se pregunta:
¿qué mide cada módulo?, ¿quién integra sus denominadores?, ¿cómo se incorpora la
selección muestral? y ¿qué conclusiones pueden sostenerse a partir de los
tabulados publicados? Estas preguntas anteceden a cualquier explicación causal.
Una regla incorrecta de codificación puede alterar la interpretación de un
porcentaje incluso si el procedimiento estadístico posterior es sofisticado.

El artículo aporta una correspondencia explícita entre preguntas, variables e
indicadores; una selección de cifras oficiales de referencia; un protocolo en
R con controles de integridad; y anexos para replicación. La sección de
resultados distingue evidencia secundaria publicada de estimaciones propias
todavía no calculadas. No se utiliza la información artificial de las pruebas
como evidencia sobre la población salvadoreña.

La presentación de Infosegura \citep{infosegura2021} introduce una segunda
pregunta: ¿en qué medida una mejora de los indicadores generales de seguridad
describe también la protección dentro del hogar? Su lectura de homicidios,
victimización, percepción y privaciones económicas permite situar los módulos
MICS en un diagnóstico más amplio. Las fuentes no observan a las mismas
personas ni el mismo momento; se articulan documentalmente, sin unir sus
registros ni construir una serie de prevalencia familiar de 2014 a 2021.

\section{Contexto y marco conceptual}
\subsection{Violencia, protección y política pública}
El informe nacional incorpora disciplina, actitudes y matrimonio temprano
dentro de su análisis de protección a la infancia. Estos módulos amplían el
diagnóstico de salud y condiciones de vida hacia relaciones de poder,
socialización y ejercicio de derechos \citep{ens}. Sin embargo, su cobertura no
equivale a la de una encuesta especializada en todas las formas de violencia.
No ofrece una medición integral de homicidios, violencia de pandillas, abuso
sexual o coerción de pareja.

El marco INSPIRE de la Organización Mundial de la Salud y sus asociados
identifica ámbitos complementarios de prevención, entre ellos normas sociales,
apoyo a cuidadores, entornos seguros, fortalecimiento económico y servicios de
respuesta \citep{inspire}. Este enfoque respalda examinar prácticas y actitudes
en conjunto, pero no justifica asumir que una respuesta normativa cause un
comportamiento individual ni que una intervención específica explique un
porcentaje transversal.

En El Salvador, las reformas sobre matrimonio infantil reconocidas por agencias
de Naciones Unidas en agosto de 2017 constituyen un antecedente relevante para
la discusión disponible a noviembre de 2021 \citep{onu2017}. La ENS MICS fue
levantada antes de esa reforma. Sus datos de 2014 pueden funcionar como
referencia previa, pero no permiten estimar el impacto de un cambio legal
posterior. Además, una unión de convivencia y un matrimonio formal son
circunstancias que el cuestionario considera dentro de una trayectoria común;
sus mecanismos sociales y jurídicos no son idénticos.

\subsection{Tres dimensiones que no deben confundirse}
La disciplina violenta constituye una experiencia reportada por un adulto
acerca de un niño durante el mes anterior. La justificación de la violencia
es una opinión declarada frente a situaciones hipotéticas. La unión temprana
es un evento de inicio de convivencia conyugal reconstruido mediante la historia declarada. Estos constructos
pueden relacionarse, pero tienen unidades, ventanas temporales y procesos de
medición diferentes.

Por ello no se suman en un índice único de violencia. Una escala conjunta
exigiría justificar ponderaciones, dirección de los componentes y comparabilidad,
además de validar su estructura. Tampoco se denomina a las preguntas de actitudes
una escala validada de machismo estructural. El término \emph{instrumento} se
refiere aquí al cuestionario de encuesta, no a una variable instrumental
econométrica.

\subsection{El hogar dentro del análisis multidimensional de seguridad}
Infosegura plantea que el control del crimen debe acompañarse de oportunidades
de desarrollo y participación ciudadana. Su agenda de prevención recomienda seguir
avisos de violencia intrafamiliar para prevenir feminicidios y actuar en los
espacios públicos de mayor uso. Esta perspectiva permite incorporar a MICS
como una fuente sobre relaciones y prácticas familiares, complementaria a
la medición de seguridad ciudadana \citep{infosegura2021}.

El vínculo analítico se establece por mecanismos que requieren evaluación:
formas de crianza, aceptación normativa del uso de fuerza, autonomía de las
adolescentes y condiciones materiales del hogar. Son hipótesis para orientar
investigación y prevención. Los agregados disponibles no permiten saber si
los niños disciplinados violentamente son quienes sufrirán o ejercerán
violencia posterior, ni si las mujeres que justifican un escenario son las
víctimas de un feminicidio registrado años después.

La articulación exige observar también el entorno. La posibilidad de usar
transporte o espacios comunitarios y el acceso al empleo pueden relacionarse
con autonomía y acceso a servicios; no se estima aquí la dirección ni magnitud
de esas relaciones. En particular, una unión iniciada antes de los 18 años
plantea preguntas sobre agencia y oportunidades, pero no identifica por sí
sola la ocurrencia de cada tipo de violencia de pareja.

\begin{table}[htbp]
\centering\small
\caption{Correspondencia entre MICS e Infosegura}\label{tab:comparabilidad}
\begin{tabularx}{\textwidth}{@{}p{2.5cm}YYp{3.4cm}@{}}\toprule
Fuente y dimensión & Población y período & Unidad & Relación con el estudio\\\midrule
MICS: disciplina & Niño seleccionado 1-14, último mes en 2014 & Porcentaje con actos declarados & Crianza y protección infantil\\
MICS: normas y unión & Mujeres 15-49; cohorte 20-24; 2014 & Opiniones y eventos previos de unión & Género y autonomía\\
Infosegura: homicidios & Población total o por sexo; 2013-2021 & Tasas por cien mil; 2021 proyectado & Contexto de violencia letal\\
Cultura de Paz & Población encuestada; 2017-2021 & Victimización en 12 meses y percepción & Experiencia delictiva y espacios públicos\\
Registros y privaciones & Cortes propios de cada gráfico & Denuncias mensuales; variaciones en pp & Demanda institucional y contexto material\\\bottomrule
\end{tabularx}
\fuente{Cuestionario MICS y PNUD/Infosegura (2021). La serie de victimización no explicita la edad del universo; no se le atribuye la base 18+ de los indicadores de uso de espacios públicos.}
\end{table}


\begin{table}[htbp]
\centering\small
\caption{Dimensiones, unidades y alcance interpretativo}\label{tab:dimensiones}
\begin{tabularx}{\textwidth}{@{}p{2.6cm}YYY@{}}\toprule
Dimensión & Unidad y período & Indicador principal & Lo que no identifica por sí sola\\\midrule
Disciplina & Niño seleccionado de 1 a 14 años; último mes & Algún acto físico o psicológico & Toda violencia sufrida a lo largo de la vida\\
Actitudes & Mujer de 15 a 49 años; opinión al entrevistar & Justificación en alguno de cinco escenarios & Victimización o conducta efectiva del esposo\\
Primera unión & Mujer de 20 a 24 años; historia de primera unión & Inicio antes de los 15 y de los 18 años & Incidencia anual o efecto causal de una reforma\\\bottomrule
\end{tabularx}
\fuente{Construcción analítica a partir del cuestionario ENS MICS. Los tres universos no se combinan como si fueran una sola muestra de individuos.}
\end{table}

\section{Fuente, diseño e instrumento}
\subsection{Organización de la encuesta}
La ENS MICS fue realizada por el Ministerio de Salud y el Instituto Nacional
de Salud, con colaboración de DIGESTYC y asistencia técnica de UNICEF.
El informe documenta un diseño multietápico por conglomerados, estratificado
por área urbana y rural dentro de los departamentos \citep{ens}. La biblioteca
de microdatos identifica la encuesta con la referencia
\var{SLV_2014_MICS_v01_M} \citep{catalogo}.

Para este estudio se utilizan conceptualmente tres archivos SPSS:
\var{hh.sav}, de hogares; \var{hl.sav}, del listado de miembros; y
\var{wm.sav}, de mujeres. El archivo de menores de cinco años no sustituye al
de disciplina infantil. Las llaves del hogar combinan \var{HH1} y \var{HH2};
la línea del miembro o de la mujer identifica a la persona dentro de ese hogar.
Antes de unir archivos se exige unicidad de estas llaves y se comprueba que
la vinculación no multiplique registros.

\begin{figure}[htbp]
\centering
\begin{tikzpicture}[node distance=0.75cm and 0.8cm,
 box/.style={draw=petroleo,rounded corners=2pt,align=center,text width=4.7cm,
 minimum height=1.05cm,font=\small},arr/.style={-{Stealth[length=2mm]},thick,draw=gris}]
\node[box] (hh) {\textbf{Hogares: hh.sav}\\Diseño, riqueza, área y selección SL};
\node[box,right=of hh] (wm) {\textbf{Mujeres: wm.sav}\\Actitudes y trayectoria de unión};
\node[box,below=of hh] (hl) {\textbf{Miembros: hl.sav}\\Línea, edad y sexo del niño};
\node[box,below=of hl] (child) {\textbf{Dominio infantil 1-14}\\Peso del hogar $\times$ SL1};
\node[box,below=2.55cm of wm] (women) {\textbf{Mujeres 15-49 / cohorte 20-24}\\Ponderador individual wmweight};
\draw[arr] (hh)--(hl);
\draw[arr] (hl)--(child);
\draw[arr] (hh)--(wm) node[midway,above,font=\scriptsize] {HH1, HH2};
\draw[arr] (wm)--(women);
\end{tikzpicture}
\caption{Arquitectura de los archivos y de los dominios analíticos}\label{fig:flujo}
\fuente{Esquema del protocolo propuesto. Compartir hogar no basta para atribuir maternidad ni para vincular automáticamente la actitud de una mujer con la disciplina de un niño.}
\end{figure}

\subsection{La presentación de Infosegura como fuente secundaria}
Se utiliza la copia facilitada de \emph{Análisis Multidimensional de Seguridad
Ciudadana. Primer Semestre 2021. El Salvador}, de 14 diapositivas. Sus pies
atribuyen las cifras a PNUD a través de Infosegura, con datos de mesas técnicas,
PNC, DIGESTYC y encuestas de Cultura de Paz. Se extrajeron agregados de las
cachés de sus gráficos y se cotejaron con una representación en PowerPoint.
Los libros Excel vinculados no están incluidos y no se intentó acceder a sus
enlaces privados. Las cachés permiten reproducir la selección utilizada,
pero no verificar ponderadores o registros primarios.

Aunque el título se refiere al primer semestre, la diapositiva 6 precisa un
corte a agosto y otras muestran proyecciones de 2021. Por ello no se impone
un período uniforme a toda la presentación. Las unidades y restricciones se
identifican por gráfico. Se excluyen del análisis cuantitativo migración y
revictimización por no disponer de bases suficientemente verificadas; tampoco
se usa la serie histórica mixta de la diapositiva 3.

\subsection{Selección del niño y disciplina}
El módulo de selección identifica un niño entre los miembros de 1 a 17 años.
\var{SL1} cuenta ese conjunto; \var{SL9B} y \var{SL9C} registran línea y edad
del seleccionado. El módulo de disciplina se aplica si tiene de 1 a 14 años;
\var{CD2} conserva su línea. La nota 5 de la página impresa 306 del informe
prescribe ajustar el peso del hogar por el número de niños de 1 a 17 años.
Este detalle es decisivo: el rango usado para seleccionar es más amplio que
el rango al que se aplica el indicador.

Las preguntas \var{CD3A}-\var{CD3K} describen prácticas usadas por la persona
informante u otros miembros del hogar. Sus respuestas sustantivas son uno
para sí y dos para no. La agresión psicológica comprende gritos e insultos.
La física incluye sacudidas, nalgadas y los distintos golpes descritos.
El subgrupo severo considera golpes en cara, cabeza u orejas y palizas fuertes
y repetidas. Golpear con un objeto, por sí solo, no se asigna automáticamente
a ese subgrupo; esto no minimiza la gravedad del acto.

\begin{table}[htbp]
\centering\small
\caption{Correspondencia de los ítems de disciplina}\label{tab:cd}
\begin{tabularx}{\textwidth}{@{}p{2.2cm}Yp{4.7cm}@{}}\toprule
Ítem & Contenido resumido & Clasificación operacional\\\midrule
CD3A, B, E & Quitar privilegios, explicar conducta y dar otra actividad & Métodos no violentos\\
CD3D, H & Gritar e insultar & Agresión psicológica\\
CD3C, F & Sacudir y dar nalgada con mano & Castigo físico\\
CD3G, J & Golpear con objeto; golpear mano, brazo o pierna & Castigo físico\\
CD3I, K & Golpear cara, cabeza u orejas; paliza repetida y fuerte & Castigo físico severo\\\bottomrule
\end{tabularx}
\fuente{Cuestionario, apéndice F, p. 474. El castigo severo es parte del físico total. Los grupos psicológico y físico pueden solaparse.}
\end{table}

Debe distinguirse el uso de \emph{algún} método no violento de \emph{solo}
disciplina no violenta. La segunda categoría exige al menos un acto A, B o E
y ausencia identificada de todos los actos violentos. Un hogar puede explicar
una conducta y también golpear: en ese caso hay algún método no violento,
pero no una práctica exclusivamente no violenta. La pregunta \var{CD4}, sobre
la necesidad de castigar físicamente, es una actitud y no se añade al indicador
de experiencia del niño.

\subsection{Diez escenarios de actitudes}
El cuestionario salvadoreño incluye diez situaciones \var{DV1A}-\var{DV1J},
con códigos uno para sí, dos para no y ocho para no sabe. Las cinco primeras
forman el compuesto comparable más habitual; las cinco adicionales amplían
la cobertura nacional. El universo abarca mujeres entrevistadas de 15 a 49
años, independientemente de su estado conyugal.

\begin{table}[htbp]
\centering\small
\caption{Escenarios de justificación de violencia contra la esposa}\label{tab:dvitems}
\begin{tabularx}{\textwidth}{@{}p{1.4cm}Yp{3.9cm}@{}}\toprule
Ítem & Situación planteada & Compuesto al que pertenece\\\midrule
DV1A & Salir sin avisarle al esposo & Cinco y diez\\
DV1B & Descuidar a los niños & Cinco y diez\\
DV1C & Discutir con el esposo & Cinco y diez\\
DV1D & Negarse a mantener relaciones sexuales & Cinco y diez\\
DV1E & Quemar la comida & Cinco y diez\\
DV1F & Hablar con otro hombre & Diez\\
DV1G & Salir mucho con amistades & Diez\\
DV1H & Hablar mucho con familiares & Diez\\
DV1I & No compartir sus ingresos & Diez\\
DV1J & Tocar las pertenencias del esposo & Diez\\\bottomrule
\end{tabularx}
\fuente{Cuestionario nacional, apéndice F, p. 507. Las descripciones son paráfrasis. La justificación declarada no mide directamente hechos de violencia sufridos o ejercidos.}
\end{table}

El compuesto de diez no sustituye al de cinco. Agregar escenarios aumenta
las oportunidades de una respuesta afirmativa y puede modificar el porcentaje
sin que cambie la disposición de las personas. Tampoco puede reconstruirse
su prevalencia sumando porcentajes marginales: una mujer puede justificar
más de una situación y las intersecciones no se conocen a partir de una tabla
de ítems separados.

\subsection{Historia de matrimonio y convivencia}
\var{MA1} distingue matrimonio actual, convivencia actual y ausencia de unión
actual. Para esta última categoría, \var{MA5} diferencia matrimonio previo
(uno), convivencia previa (dos) y nunca unión (tres). Confundir el código dos
con nunca unida clasificaría incorrectamente a antiguas convivientes.
\var{MA9} registra edad al inicio de la primera unión; \var{WB2}, edad actual.

La implementación utiliza edades declaradas y las contrasta con la trayectoria
conyugal. El archivo también dispone de edades derivadas \var{wage} y
\var{wagem}. Una réplica exacta de los tabulados debe comparar estas variables
y documentar las reglas de edición e imputación de fechas. No se sustituyen
silenciosamente unas por otras. Si consta unión previa y falta edad, el evento
se conserva como desconocido. Las mujeres nunca unidas integran el denominador
con resultado cero.

El dominio principal para ambos umbrales es 20 a 24 años. A diferencia de una
adolescente de 16 años, todas sus integrantes han superado el umbral de 18.
Esta restricción evita exposición incompleta respecto al evento, aunque no
elimina sesgos de recuerdo, supervivencia o migración. Una edad exactamente
igual a 15 no representa una unión \emph{antes} de los 15.

\FloatBarrier
\section{Estrategia estadística y reproducibilidad}
\subsection{Universos, entrevistas y llaves}
El ejecutor exige entrevistas completas de hogar y de mujer, según
\var{HH9=1} y \var{WM7=1}. Comprueba unicidad de hogares, miembros y mujeres;
coincidencia entre el conteo de miembros de 1 a 17 y \var{SL1}; identificación
del seleccionado en el listado; y coherencia entre las líneas \var{SL9B} y
\var{CD2} para el dominio infantil. Un diseño incompleto o un peso no positivo
detiene el análisis, en lugar de eliminar observaciones silenciosamente.

La unión de archivos selecciona campos explícitamente para evitar sufijos
ambiguos cuando mujeres ya incorpora variables del hogar. Se conserva la
distinción entre el ponderador individual y el del hogar. Las pruebas incluyen
hogares con más de un niño elegible y más de una mujer, casos que permiten
detectar multiplicación de filas y uso del ponderador incorrecto.

\subsection{Ponderación y dominios}
Sea $w_h$ el peso normalizado del hogar y $N_h^{1:17}$ su número de miembros
de 1 a 17 años. Para el niño seleccionado se usa
\begin{equation}
 w_{hi}^{\mathrm{disc}}=w_h N_h^{1:17}.
 \label{eq:peso}
\end{equation}
La multiplicación compensa la selección intrahogar documentada. No debe
reemplazarse $N_h^{1:17}$ por el número de niños de 1 a 14 o menores de cinco.
Para mujeres se utiliza \var{wmweight}, sin sustitución por \var{hhweight}.

Los diseños se construyen antes de restringir los dominios. En infancia se
retienen todos los hogares completos, incluidos aquellos cuyo seleccionado
tiene 15 a 17 años. Los hogares sin niños reciben un peso auxiliar positivo
para construir el diseño, pero contribución nula al dominio y resultado
desconocido, no cero, en disciplina. En mujeres, la cohorte 20 a 24 se obtiene
como subpoblación del diseño completo de entrevistadas.

Se emplean \var{psu} y \var{stratum} distribuidos en el archivo de hogares.
No se sustituye estrato por región. La aproximación de conglomerado último
usa linealización de Taylor y no incorpora corrección por población finita,
al no disponer de fracciones de muestreo. La estructura y la variación de
pesos son parte de la información que debe revisarse \citep{lumley}.

\subsection{Construcción con respuestas parciales}
Para un compuesto $Y_i$ basado en ítems binarios, una respuesta afirmativa
identifica presencia aunque otros componentes sean desconocidos. La ausencia
solo se identifica si todos son negativos. Si no se cumple ninguna condición,
$Y_i$ permanece desconocido. Esta regla impide que el uso automático de
sumas con eliminación de ausentes convierta patrones sin información en ceros.

Sea $D$ el dominio y $O\subseteq D$ el conjunto clasificable para un indicador.
La razón ponderada propuesta es
\begin{equation}
 \widehat p_O=\frac{\sum_{i\in O}w_iY_i}{\sum_{i\in O}w_i}.
 \label{eq:prevalencia}
\end{equation}
Cada indicador puede tener su propio conjunto $O$. Extrapolar esta razón a
todo $D$ requiere supuestos sobre la falta de respuesta. El programa exporta
tamaños sin ponderar y fracción ponderada desconocida, y no impone una muestra
completa común a variables que no la necesitan.

Para visualizar la incertidumbre por información faltante se calculan los
límites extremos
\begin{align}
 L_D&=\frac{\sum_{i\in O}w_iY_i}{\sum_{i\in D}w_i},\\
 U_D&=\frac{\sum_{i\in O}w_iY_i+\sum_{i\in D\setminus O}w_i}
 {\sum_{i\in D}w_i}.
 \label{eq:limites}
\end{align}
La diferencia $U_D-L_D$ es la fracción ponderada desconocida. Estos límites
no son intervalos de confianza; no incorporan variación muestral ni corrigen
errores de declaración. Su función es mostrar cuánto puede importar la parte
no clasificable bajo asignaciones extremas.

\subsection{Precisión, desagregación y límites de inferencia}
Los errores estándar provienen de \var{svymean}; los IC95 se calculan con
\var{svyciprop}, método logit y grados de libertad del diseño del dominio
clasificable. Se declara el ajuste de estratos con una UPM y se audita su
presencia. Dominios con menos de 25 observaciones clasificables no presentan
porcentaje; proporciones exactamente cero o uno conservan su punto, pero no
reciben un IC logit fabricado.

Las desagregaciones previstas son residencia, quintil y edad, más sexo en
disciplina. La unión temprana se desagrega por residencia y quintil dentro
de 20 a 24 años. Las categorías desconocidas de covariables se conservan
explícitamente. Una diferencia entre grupos no constituye automáticamente
una diferencia estadísticamente respaldada ni una relación causal.

No se aplica emparejamiento por puntaje de propensión. No se ha definido un
tratamiento con temporalidad y supuestos de identificación. La riqueza
observada al entrevistar podría ser antecedente o consecuencia de una unión
iniciada años antes. El ajuste por covariables contemporáneas no resuelve esa
ambigüedad. Tampoco se atribuye maternidad a una mujer solo por compartir hogar
con un niño seleccionado.

\subsection{Triangulación documental y proceso de registro}
La integración compara dimensiones y patrones de fuentes independientes.
No agrega porcentajes de niños, tasas de mortalidad y opiniones de adultos
en una misma escala. Tampoco asigna el contexto económico de 2020 a cada
niño observado en 2014 como si fuese una covariable contemporánea.
El cuadro \ref{tab:comparabilidad} fija la correspondencia entre fuentes.

En los registros de denuncias, la observación depende tanto de la ocurrencia
como del acceso al sistema y su captura. Una descomposición conceptual es
\begin{equation}
 \mathbb{E}[R_t] = V_t\,\pi_t\,\rho_t,
 \label{eq:registro}
\end{equation}
donde $V_t$ representa hechos potencialmente denunciables, $\pi_t$ la
probabilidad de denunciar y $\rho_t$ la de registro una vez denunciado el
hecho. Es un esquema de medición, no un modelo estimado con estas fuentes.
Un cambio de denuncias no identifica por sí solo un cambio equivalente de
violencia real, especialmente si cambian movilidad, atención o período cubierto.

\section{Evidencia publicada y resultados de la auditoría}
\subsection{Naturaleza de las cifras presentadas}
Los cuadros y gráficos siguientes transcriben una selección de cifras del
informe oficial. Constituyen \textbf{evidencia secundaria publicada}; no son
resultados de ejecutar el programa de este repositorio sobre los microdatos.
Los valores se conservaron a un decimal y se cotejaron visualmente con las
tablas CP.5, CP.7 y CP.10. No se reconstruyen errores estándar o intervalos
a partir de porcentajes redondeados.

El informe contiene algunas inconsistencias editoriales: el texto introductorio
de CP.5 se refiere a trabajo infantil aunque su encabezado y columnas son de
disciplina; además, ciertas filas departamentales de CP.10 muestran valores
que no son coherentes con sus encabezados. Esta edición emplea únicamente
celdas nacionales, de sexo, área o cohorte visualmente identificables; no
reproduce las filas departamentales problemáticas. Una réplica con microdatos
deberá aclarar esas discrepancias antes de ampliar la comparación.

\subsection{Disciplina infantil}
La fila nacional de CP.5 informa 52,0\,\% para cualquier disciplina violenta,
32,2\,\% para agresión psicológica, 39,4\,\% para castigo físico y 3,0\,\% para
castigo severo. La categoría exclusivamente no violenta alcanza 35,5\,\%.
Estas categorías no deben sumarse: las dos formas de violencia pueden
coexistir y el castigo severo es un subconjunto del físico.

\begin{figure}[htbp]
\centering
\begin{tikzpicture}
\begin{axis}[width=0.86\textwidth,height=6.3cm,xbar,bar width=11pt,
 xmin=0,xmax=60,xtick={0,10,20,30,40,50,60},
 ytick={0,1,2,3,4},yticklabels={Solo no violenta,Psicológica,Física total,Física severa,Cualquier violenta},
 ymin=-0.7,ymax=4.7,xlabel={Porcentaje de niños de 1 a 14 años},
 axis x line*=bottom,axis y line*=left,xmajorgrids=true,grid style={gray!15},
 tick label style={font=\small},nodes near coords,
 nodes near coords style={font=\small,/pgf/number format/fixed,/pgf/number format/precision=1}]
\addplot[fill=petroleo,draw=none] table[x=porcentaje,y=pos,col sep=comma]{data/disciplina.csv};
\end{axis}
\end{tikzpicture}
\caption{Métodos de disciplina: cifras nacionales publicadas}\label{fig:disciplina}
\fuente{ENS MICS, tabla CP.5, p. 282. Transcripción de porcentajes publicados, sin IC propios. Categorías solapadas salvo ``solo no violenta'' respecto a violencia.}
\end{figure}

\begin{table}[htbp]
\centering\small
\caption{Disciplina infantil por sexo y residencia: porcentajes publicados}\label{tab:discgrupo}
\begin{tabularx}{\textwidth}{@{}Yrrrrr@{}}\toprule
Grupo & Solo no & Psico- & Físico & Físico & Cualquier\\
 & violenta & lógica & total & severo & violenta\\\midrule
Total & 35,5 & 32,2 & 39,4 & 3,0 & 52,0\\
Masculino & 34,3 & 33,2 & 42,6 & 3,8 & 54,6\\
Femenino & 36,6 & 31,3 & 36,2 & 2,2 & 49,5\\
Urbana & 35,6 & 34,2 & 40,5 & 3,0 & 53,8\\
Rural & 35,4 & 29,6 & 37,9 & 3,0 & 49,6\\\bottomrule
\end{tabularx}
\fuente{Informe oficial, CP.5, p. 282. No se interpretan diferencias puntuales como significativas sin la varianza del diseño.}
\end{table}

La diferencia publicada entre niños y niñas en cualquier disciplina violenta
es de 5,1 puntos porcentuales; la diferencia urbana-rural es de 4,2 puntos.
Son restas descriptivas de celdas redondeadas, no contrastes inferenciales.
Estos datos tampoco autorizan a concluir que el área de residencia produzca
la violencia: composición de los hogares, declaración y otros factores
podrían contribuir a la diferencia.

\subsection{Justificación declarada de la violencia}
La tabla CP.10 informa 7,7\,\% de justificación en al menos uno de los primeros
cinco escenarios. El valor publicado es 5,9\,\% en área urbana y 10,8\,\% en
área rural. Entre los ítems individuales, descuidar a los niños presenta
5,2\,\%, mientras que hablar con otro hombre, perteneciente a la ampliación
nacional, registra 5,6\,\%.

\begin{figure}[htbp]
\centering
\begin{tikzpicture}
\begin{axis}[width=0.80\textwidth,height=8.1cm,xbar,bar width=8pt,
 xmin=0,xmax=7,xtick={0,1,2,3,4,5,6,7},
 ytick={0,1,2,3,4,5,6,7,8,9},
 yticklabels={Salir sin avisar,Descuidar niños,Discutir,Negarse a sexo,Quemar comida,
 Hablar con otro hombre,Salir con amistades,Hablar con familiares,No compartir ingresos,Tocar pertenencias},
 ymin=-0.7,ymax=9.7,xlabel={Porcentaje de mujeres de 15 a 49 años},
 axis x line*=bottom,axis y line*=left,xmajorgrids=true,grid style={gray!15},
 tick label style={font=\small},nodes near coords,
 nodes near coords style={font=\small,/pgf/number format/fixed,/pgf/number format/precision=1}]
\addplot[fill=petroleo,draw=none] table[x=porcentaje,y=pos,col sep=comma]{data/actitudes.csv};
\end{axis}
\end{tikzpicture}
\caption{Justificación por escenario: porcentajes nacionales publicados}\label{fig:actitudes}
\fuente{ENS MICS, CP.10, p. 296, fila Total. A-E forman el compuesto de cinco; F-J son adicionales. No sumar las barras ni interpretarlas como victimización.}
\end{figure}

\begin{table}[htbp]
\centering\small
\caption{Justificación de cinco escenarios y contraste de universos}\label{tab:actgrupos}
\begin{tabularx}{\textwidth}{@{}Yrp{6.8cm}@{}}\toprule
Dominio & Porcentaje & Interpretación\\\midrule
Mujeres 15-49, total & 7,7 & Algún escenario A-E, según tabulado oficial\\
Área urbana, 15-49 & 5,9 & Mismo conjunto A-E y grupo de edades\\
Área rural, 15-49 & 10,8 & Diferencia rural-urbana descriptiva: 4,9 pp\\
Mujeres 20-24 & 6,3 & Subgrupo de edad, no todas las mujeres\\
Compuesto de diez & \multicolumn{1}{c}{---} & No recuperable por suma de marginales; estimación propia pendiente\\\bottomrule
\end{tabularx}
\fuente{CP.10, p. 296. El guion indica ausencia de una cifra verificada para el compuesto de diez en esta transcripción; no equivale a cero.}
\end{table}

El contraste rural-urbano difiere del patrón puntual observado en disciplina.
Esto muestra por qué las dimensiones no pueden tratarse como sustitutas.
No demuestra una relación individual negativa entre actitudes y prácticas:
los informantes, las unidades y los universos son diferentes. Inferir una
relación individual de estos agregados incurriría en una lectura ecológica
sin respaldo de vinculación personal.

\subsection{Primera unión: importancia de la cohorte}
La fila de 20 a 24 años de CP.7 publica 5,8\,\% de unión antes de los 15 y
25,5\,\% antes de los 18. En cambio, el total de 20 a 49 años alcanza 7,3\,\%
y 29,1\,\%, respectivamente. Son dominios distintos. La diferencia entre ellos
no constituye una tendencia temporal estimada: mezcla cohortes e historias
de exposición y selección.

\begin{figure}[htbp]
\centering
\begin{tikzpicture}
\begin{axis}[width=0.86\textwidth,height=5.7cm,ybar,bar width=30pt,
 ymin=0,ymax=35,ytick={0,5,10,15,20,25,30,35},
 xmin=13,xmax=20,xtick={15,18},xticklabels={Antes de 15,Antes de 18},
 ylabel={Porcentaje de mujeres de 20 a 24 años},
 axis x line*=bottom,axis y line*=left,ymajorgrids=true,grid style={gray!15},
 nodes near coords,nodes near coords style={font=\small},tick label style={font=\small}]
\addplot[fill=ocre,draw=none] table[x=umbral,y=porcentaje,col sep=comma]{data/union.csv};
\end{axis}
\end{tikzpicture}
\caption{Primera unión en mujeres de 20 a 24 años: cifras publicadas}\label{fig:union}
\fuente{ENS MICS, CP.7, p. 288, fila 20-24. El evento antes de 15 está contenido en el evento antes de 18; las barras no son categorías excluyentes.}
\end{figure}

\begin{table}[htbp]
\centering\small
\caption{Umbral de primera unión y población de referencia}\label{tab:union}
\begin{tabularx}{\textwidth}{@{}YrrY@{}}\toprule
Población & Antes de 15 & Antes de 18 & Uso en este estudio\\\midrule
Mujeres 15-49 & 7,0 & --- & Referencia del informe; no es la cohorte principal\\
Mujeres 20-49 & 7,3 & 29,1 & Referencia amplia, distinta de 20-24\\
Mujeres 20-24 & 5,8 & 25,5 & Cohorte principal para ambos umbrales\\\bottomrule
\end{tabularx}
\fuente{CP.7, pp. 287-288. Valores porcentuales publicados. No se calcula el evento antes de 18 sobre todo el grupo 15-49 debido a exposición incompleta en parte de ese universo.}
\end{table}

La inclusión de nunca unidas en el denominador es esencial. Si solo se
consideraran mujeres que alguna vez convivieron en unión, el parámetro sería
la proporción de inicios tempranos \emph{entre quienes se unieron}, que no es
la prevalencia de unión temprana en toda la cohorte. El programa conserva esa
distinción y comprueba coherencia con edad actual.

\subsection{Violencia letal: mejora general y diferencias por género}
La serie de Infosegura registra una tasa total de homicidios de 106,3
por cien mil habitantes en 2015 y de 20,9 en 2020. La proyección rotulada para
2021 es 20,8. En mujeres, la tasa correspondiente es 17,3 en 2015, 3,8 en
2020 y 4,9 en la proyección de 2021. La serie de violencia de género informa tasas de
feminicidio de 2,2 en 2020 y 3,1 en su proyección de 2021. Se preserva la
condición proyectada del último punto, sin presentarlo como cierre anual.

\begin{figure}[htbp]
\centering
\begin{minipage}[t]{0.49\textwidth}\vspace{0pt}\centering
\begin{tikzpicture}
\begin{axis}[width=0.9\linewidth,height=5.7cm,title={Homicidios: total},
 xmin=2013,xmax=2021,ymin=0,ymax=120,xtick={2013,2015,2017,2019,2021},
 xticklabel style={rotate=45,anchor=east},scaled x ticks=false,
 /pgf/number format/1000 sep={},ylabel={Por cien mil habitantes},
 tick label style={font=\scriptsize},title style={font=\small},
 grid=major,grid style={gray!15},filter discard warning=false]
\addplot[petroleo,thick,mark=*,restrict x to domain=2013:2020]
 table[x=anio,y=total,col sep=comma]{data/infosegura_homicidios.csv};
\addplot[petroleo,thick,dashed,mark=o,restrict x to domain=2020:2021]
 table[x=anio,y=total,col sep=comma]{data/infosegura_homicidios.csv};
\end{axis}\end{tikzpicture}
\end{minipage}\hfill
\begin{minipage}[t]{0.49\textwidth}\vspace{0pt}\centering
\begin{tikzpicture}
\begin{axis}[width=0.9\linewidth,height=5.7cm,title={Muertes de mujeres},
 xmin=2013,xmax=2021,ymin=0,ymax=20,xtick={2013,2015,2017,2019,2021},
 xticklabel style={rotate=45,anchor=east},scaled x ticks=false,
 /pgf/number format/1000 sep={},ylabel={Por cien mil mujeres},
 tick label style={font=\scriptsize},title style={font=\small},
 grid=major,grid style={gray!15},filter discard warning=false,
 legend style={font=\scriptsize,at={(0.5,-0.45)},anchor=north,draw=none}]
\addplot[petroleo,thick,mark=*,restrict x to domain=2013:2020]
 table[x=anio,y=mujeres,col sep=comma]{data/infosegura_homicidios.csv};
\addlegendentry{Homicidios de mujeres}
\addplot[ocre,thick,mark=square*,restrict x to domain=2013:2020]
 table[x=anio,y=feminicidios,col sep=comma]{data/infosegura_homicidios.csv};
\addlegendentry{Feminicidios}
\addplot[petroleo,thick,dashed,mark=o,forget plot,restrict x to domain=2020:2021]
 table[x=anio,y=mujeres,col sep=comma]{data/infosegura_homicidios.csv};
\addplot[ocre,thick,dashed,mark=square,forget plot,restrict x to domain=2020:2021]
 table[x=anio,y=feminicidios,col sep=comma]{data/infosegura_homicidios.csv};
\end{axis}\end{tikzpicture}
\end{minipage}
\caption{Violencia letal en Infosegura: tasas y proyecciones, 2013-2021}\label{fig:infosegura_letal}
\fuente{PNUD/Infosegura (2021), con datos de mesas técnicas y DIGESTYC. El tramo discontinuo termina en la \textbf{proyección de 2021}. Los paneles usan escalas y poblaciones distintas. El total de homicidios de mujeres incluye feminicidios; las dos series del panel derecho no se suman.}
\end{figure}

El contraste es relevante para el argumento del artículo: una mejora de la
tasa nacional no permite asumir una evolución homogénea por género. MICS
agrega información sobre prácticas y normas que los registros letales no
describen. La coexistencia documental de disciplina violenta en 2014 y
reducción de homicidios después de 2015 no demuestra continuidad individual,
persistencia cuantificada de violencia familiar en 2021 ni efectos de una
política específica.

El cuadro \ref{tab:cambio_letal} cuantifica el descenso entre dos años
observados. La tasa total cae 80,3\,\%, pero la de feminicidios lo hace
72,3\,\%, aproximadamente ocho puntos menos en términos de reducción
relativa. Se trata de porcentajes de cambio de tasas, distintos de puntos
porcentuales de prevalencia. Las bases iniciales son diferentes; una menor
caída relativa no implica mayor número absoluto de hechos. La clasificación
de feminicidio exige además considerar prácticas de registro y tipificación.

\begin{table}[htbp]\centering\small
\caption{Cambio observado de tasas de violencia letal, 2015--2020}\label{tab:cambio_letal}
\begin{tabularx}{\textwidth}{@{}Yrrrr@{}}\toprule
Serie & 2015 & 2020 & Diferencia & Cambio (\%)\\\midrule
Homicidios, total & 106,3 & 20,9 & -85,4 & -80,3\\
Homicidios, hombres & 206,7 & 39,7 & -167,0 & -80,8\\
Homicidios, mujeres & 17,3 & 3,8 & -13,4 & -77,8\\
Feminicidios & 7,8 & 2,2 & -5,7 & -72,3\\
\bottomrule

\end{tabularx}
\fuente{Cálculos del artículo sobre las series de PNUD/Infosegura (2021), con precisión disponible antes del redondeo a una décima. Tasas por cien mil habitantes del grupo correspondiente. Diferencia en unidades de tasa; cambio: $100(r_{2020}/r_{2015}-1)$. No interviene la proyección de 2021.}
\end{table}

Los agregados de Infosegura también incluyen promedios mensuales de denuncias
intrafamiliares de 96,1, 77,9 y 105,4 para 2019, 2020 y 2021, respectivamente.
La fuente no precisa el corte de 2021. Se conservan como promedios reportados,
sin anualizarlos ni calcular con ellos una variación de
incidencia familiar. La ecuación \ref{eq:registro} ayuda a explicar por qué
una denuncia no equivale a una medición exhaustiva de violencia ocurrida.

La nota institucional de octubre de 2021 mantiene una proyección nacional
de 20,8 homicidios por cien mil habitantes y advierte una reducción desigual
por sexo \citep{pnud2021octubre}. Esa lectura respalda desagregar las
mejoras; no convierte una proyección en cierre anual ni prueba un mecanismo
familiar. Sus cifras resumidas de denuncias no se interpretan como totales
anuales comparables: aquí se conservan las unidades de promedio mensual.

\subsection{Victimización por delincuencia y exposición territorial}
La serie de victimización en los últimos doce meses registra 14,0\,\%
en 2017, 13,9\,\% en 2018, 10,8\,\% en 2019 y 10,4\,\% en 2021. En este
último año los valores son 10,7\,\% para hombres y 10,2\,\% para mujeres,
mientras área urbana y rural registran 12,7\,\% y 6,8\,\%. Se conservan las
definiciones de la fuente: estas cifras de delincuencia no sustituyen
la prevalencia de disciplina o violencia de pareja.

\begin{figure}[htbp]
\centering
\begin{minipage}[t]{0.49\textwidth}\vspace{0pt}\centering
\begin{tikzpicture}
\begin{axis}[width=0.9\linewidth,height=5.4cm,title={Por sexo},
 xmin=2017,xmax=2021,ymin=0,ymax=20,xtick={2017,2018,2019,2021},
 scaled x ticks=false,/pgf/number format/1000 sep={},ylabel={Porcentaje},
 grid=major,grid style={gray!15},tick label style={font=\scriptsize},
 title style={font=\small},legend style={font=\scriptsize,at={(0.5,-0.25)},anchor=north,draw=none}]
\addplot[petroleo,thick,mark=*] table[x=anio,y=hombres,col sep=comma]{data/infosegura_victimizacion.csv};
\addlegendentry{Hombres}
\addplot[ocre,thick,mark=square*] table[x=anio,y=mujeres,col sep=comma]{data/infosegura_victimizacion.csv};
\addlegendentry{Mujeres}
\end{axis}\end{tikzpicture}
\end{minipage}\hfill
\begin{minipage}[t]{0.49\textwidth}\vspace{0pt}\centering
\begin{tikzpicture}
\begin{axis}[width=0.9\linewidth,height=5.4cm,title={Por residencia},
 xmin=2017,xmax=2021,ymin=0,ymax=20,xtick={2017,2018,2019,2021},
 scaled x ticks=false,/pgf/number format/1000 sep={},ylabel={Porcentaje},
 grid=major,grid style={gray!15},tick label style={font=\scriptsize},
 title style={font=\small},legend style={font=\scriptsize,at={(0.5,-0.25)},anchor=north,draw=none}]
\addplot[petroleo,thick,mark=*] table[x=anio,y=urbana,col sep=comma]{data/infosegura_victimizacion.csv};
\addlegendentry{Urbana}
\addplot[ocre,thick,mark=square*] table[x=anio,y=rural,col sep=comma]{data/infosegura_victimizacion.csv};
\addlegendentry{Rural}
\end{axis}\end{tikzpicture}
\end{minipage}
\caption{Victimización por delincuencia en los últimos doce meses}\label{fig:infosegura_vict}
\fuente{PNUD/Infosegura (2021), con encuestas de Cultura de Paz. Años observados: 2017, 2018, 2019 y 2021; no se imputa 2020. Los agregados no explicitan la edad del universo ni incluyen IC; las diferencias no se presentan como significativas.}
\end{figure}

La brecha entre hombres y mujeres pasa de 5,39 puntos en 2018 a 1,39 en
2019 y 0,49 en 2021. La urbana-rural, en cambio, pasa de 2,43 en 2019 a
5,86 en 2021; esta última distancia se aproxima a los 6,37 puntos de 2018.
Por tanto, la caída del total no describe por sí sola la distribución
territorial. Al no disponer de errores estándar y covarianzas entre
dominios, la reducción de una brecha y ampliación de otra son contrastes
puntuales, no evidencia de cambios estadísticamente significativos. Un
diagnóstico de cobertura debe conservar ambos ejes de desagregación.

La brecha urbana-rural de victimización tiene el mismo signo que la diferencia
descriptiva de disciplina de MICS y distinto signo que la justificación de
violencia contra la esposa. Esa observación no es una correlación estimada:
son años, unidades y experiencias diferentes. Su utilidad consiste en evitar
que un indicador territorial se convierta en sustituto automático de todos
los demás y en orientar una agenda de desagregación comparable.

\subsection{Condiciones materiales, prioridades y percepción}
Los indicadores de opinión sitúan economía y seguridad como principal problema del país
en 46,9\,\% y 26,6\,\% en 2021. En 2019 informa 32,3\,\% y 59,8\,\%.
La serie incorpora además COVID-19 como categoría en 2021, por lo que el
cambio de composición debe reconocerse. No se interpreta como una medida
de ingresos, desempleo o victimización individual.

\begin{table}[htbp]
\centering\small
\caption{Indicadores de contexto económico y seguridad ciudadana}\label{tab:infosegura_contexto}
\begin{tabularx}{\textwidth}{@{}Yp{2.5cm}p{2.1cm}Y@{}}\toprule
Indicador & Período y unidad & Valor & Lectura para el estudio\\\midrule
Economía, principal problema & 2021; \% & 46,9 & Prioridad declarada, no ingreso observado\\
Seguridad, principal problema & 2021; \% & 26,6 & No es prevalencia de delitos\\
Privación por desempleo & 2019-2020; cambio en pp & +10,4 & Variación de un indicador de privación, no tasa de desempleo\\
Restricciones por inseguridad & 2019-2020; cambio en pp & $-6,8$ & Componente de privaciones; distinto de disciplina familiar\\
Inseguridad en transporte & 2021; \% reportado & 68,1 & Base por espacio no detallada en la fuente\\
Inseguridad en comunidad & 2021; \% reportado & 19,9 & Percepción, no hecho delictivo registrado\\\bottomrule
\end{tabularx}
\fuente{PNUD/Infosegura (2021). La base 18+ del indicador de uso de espacios no se atribuye automáticamente a los indicadores de percepción. Las unidades de las filas no se suman.}
\end{table}

El aumento del componente de privación por desempleo plantea preguntas sobre
recursos para la crianza, dependencia económica y acceso a servicios. Estas
relaciones son hipótesis; la variación de 10,4 puntos no demuestra que el
desempleo causara las prácticas observadas por MICS siete años antes.
La agenda de prevención de Infosegura también propone atender los espacios públicos de mayor
uso. Su aportación es ampliar la prevención hacia condiciones de acceso y
movilidad, junto con la protección dentro del hogar.

\subsubsection{Desplazamiento de prioridades y persistencia del temor}
Entre 2019 y 2021, la economía como principal problema aumenta 14,7 puntos,
mientras seguridad disminuye 33,2 puntos (figura \ref{fig:prioridades}). En
2021, COVID-19 representa 22,0\,\%. La incorporación de esa categoría cambia
el conjunto de respuestas: las variaciones de seguridad y economía no son
una redistribución exclusiva entre ambas. Tampoco miden la intensidad de
cada preocupación, pues elegir un problema principal puede coexistir con
preocupaciones importantes en otros ámbitos.

\begin{figure}[htbp]\centering
\begin{tikzpicture}\begin{axis}[width=0.85\textwidth,height=5.7cm,
 xmin=2016.8,xmax=2021.2,ymin=0,ymax=80,xtick={2017,2018,2019,2021},
 scaled x ticks=false,/pgf/number format/1000 sep={},ylabel={Porcentaje},
 grid=major,grid style={gray!15},tick label style={font=\small},
 legend style={at={(0.5,-0.18)},anchor=north,legend columns=3,column sep=6pt,draw=none,font=\small}]
\addplot[petroleo,thick,mark=*] table[x=anio,y=seguridad,col sep=comma]{data/infosegura_preocupaciones.csv};
\addlegendentry{Seguridad}
\addplot[ocre,thick,mark=square*] table[x=anio,y=economia,col sep=comma]{data/infosegura_preocupaciones.csv};
\addlegendentry{Economía}
\addplot[gris,only marks,mark=triangle*,mark size=3pt] coordinates {(2021,22.02)};
\addlegendentry{COVID-19}
\end{axis}\end{tikzpicture}
\caption{Evolución de las prioridades declaradas: principal problema del país}\label{fig:prioridades}
\fuente{PNUD/Infosegura (2021). Series seleccionadas, no distribución completa de respuestas. COVID-19 solo se observa como categoría en 2021; los años anteriores no se codifican como cero. Sin dato de 2020 ni IC propios.}
\end{figure}

La percepción de inseguridad también disminuye, pero conserva diferencias
marcadas por espacio (figura \ref{fig:percepcion}). Entre 2019 y 2021 cae
10,7 puntos en transporte público y 12,5 en comunidad. En 2021, las cifras
reportadas son 68,1\,\% y 19,9\,\%, respectivamente. Su distancia descriptiva
de 48,2 puntos muestra que un mismo diagnóstico nacional requiere observar
lugares de movilidad y de residencia. La fuente no detalla las bases de
cada espacio; la distancia no se interpreta como una comparación de riesgo
individual ajustada por exposición.

\begin{figure}[htbp]\centering
\begin{tikzpicture}\begin{axis}[width=0.85\textwidth,height=5.7cm,
 xmin=2016.8,xmax=2021.2,ymin=0,ymax=100,xtick={2017,2018,2019,2021},
 scaled x ticks=false,/pgf/number format/1000 sep={},ylabel={Porcentaje reportado},
 grid=major,grid style={gray!15},tick label style={font=\small},
 legend style={at={(0.5,-0.18)},anchor=north,legend columns=2,column sep=6pt,draw=none,font=\small}]
\addplot[petroleo,thick,mark=*] table[x=anio,y=transporte_publico,col sep=comma]{data/infosegura_percepcion.csv};
\addlegendentry{Transporte público}
\addplot[ocre,thick,mark=square*] table[x=anio,y=comunidad,col sep=comma]{data/infosegura_percepcion.csv};
\addlegendentry{Comunidad}
\end{axis}\end{tikzpicture}
\caption{Percepción de inseguridad en transporte y comunidad}\label{fig:percepcion}
\fuente{PNUD/Infosegura (2021). La base de personas de 18+ del indicador de uso de espacios no se traslada a estas series de percepción. Percepción distinta de victimización; segmentos entre años observados, sin imputar 2020.}
\end{figure}

La combinación de menor victimización declarada y persistencia del temor en
transporte sugiere evaluar varias dimensiones de seguridad a la vez. Para
las familias, el acceso a escuela, empleo y servicios requiere movilidad;
el temor puede ser relevante incluso sin una experiencia delictiva reciente.
Esta es una interpretación para orientar preguntas de investigación, no una
estimación del efecto del transporte sobre disciplina o primera unión. Se
necesitaría medir uso efectivo, frecuencia de desplazamiento y horarios para
examinar esa relación.

\subsubsection{Privaciones: cambios heterogéneos entre 2019 y 2020}
La figura \ref{fig:privaciones} examina los veinte cambios de privación disponibles.
Siete indicadores aumentan, once disminuyen y dos quedan sin cambio en la
precisión publicada. El aumento de privación por desempleo (+10,4 puntos)
destaca frente a otras variaciones positivas de hasta 1,3 puntos. Al mismo
tiempo, disminuyen restricciones por inseguridad ($-6{,}8$), falta de acceso
a salud ($-3{,}0$) e incidencia de crimen y delito ($-2{,}4$). El conjunto
documenta movimientos de distinto signo, relevantes para una interpretación
multidimensional del bienestar.

\begin{figure}[!htbp]\centering
\begin{tikzpicture}\begin{axis}[width=0.76\textwidth,height=11.0cm,xbar,
 bar width=7pt,xmin=-8,xmax=12,xtick={-8,-4,0,4,8,12},
 ymin=-0.7,ymax=19.7,ytick={0,...,19},
 yticklabels={Restricciones por inseguridad,Sin acceso a salud,Crimen y delito,
 Hacinamiento,Tenencia del terreno,Cuidado temprano inadecuado,
 Sin seguridad social,Sin saneamiento,Inasistencia escolar,
 Subempleo e inestabilidad,Trabajo infantil,Rezago educativo,
 Sin agua potable,Pared y piso inadecuados,Sin espacios de esparcimiento,
 Baja educación de adultos,Techo inadecuado,Inseguridad alimentaria,
 Daños y riesgos ambientales,Desempleo},
 xlabel={Variación 2019--2020 (puntos porcentuales)},
 axis x line*=bottom,axis y line*=left,xmajorgrids=true,grid style={gray!15},
 tick label style={font=\scriptsize},nodes near coords,
 nodes near coords style={font=\scriptsize,/pgf/number format/fixed,/pgf/number format/precision=1}]
\addplot[fill=petroleo,draw=none] table[x=cambio,y=pos,col sep=comma]{data/derivada_privaciones.csv};
\draw[gris,thick] (axis cs:0,-0.7)--(axis cs:0,19.7);
\end{axis}\end{tikzpicture}
\caption{Variación de los indicadores de privación}\label{fig:privaciones}
\fuente{PNUD/Infosegura (2021). Etiquetas abreviadas; textos originales en CSV. Cambios de indicadores de privación, no sus niveles. El componente desempleo no es la tasa de desempleo. Cero es ausencia de cambio al redondeo disponible.}
\end{figure}

La magnitud de desempleo invita a examinar apoyos de ingreso y cuidado junto
con servicios de protección. Sin embargo, no se suman las veinte variaciones
para obtener una variación de pobreza multidimensional: las privaciones
pueden coexistir en un hogar y el índice depende de ponderaciones, umbral y
distribución conjunta. Una reducción del indicador de inseguridad tampoco
implica que el nivel de privación sea bajo o que haya desaparecido. Sin
niveles iniciales, no pueden compararse niveles finales solo por la altura
de estas barras.

La lectura conjunta con MICS permite formular una agenda concreta: distinguir
estrés material, aceptación de prácticas violentas y exposición efectiva;
observar sus intersecciones con datos compatibles; y evaluar si los servicios
alcanzan a hogares con varias necesidades. Los agregados de 2019--2021 no
explican por sí solos la conducta familiar de 2014. Su contribución al
artículo consiste en precisar las dimensiones que una siguiente medición
debe integrar, manteniendo sus distintos períodos y poblaciones.



\subsection{Hallazgos de la auditoría y ejecución pendiente}
La auditoría documental identificó correcciones en la propuesta previa:
minúsculas coherentes con la limpieza de nombres; selección intrahogar por
miembros 1 a 17; código correcto de nunca unión; separación de ``algún'' y
``solo'' método no violento; y construcción de dominios después del diseño.
Su finalidad es eliminar errores de definición, no garantizar de antemano
coincidencia con cifras oficiales.

El programa se verificó con seis bloques y 35 comprobaciones automatizadas,
incluida la lectura y escritura de archivos SPSS artificiales y la exportación
de tablas y gráficos. Se restauraron 72 paquetes fijados en el archivo de
dependencias. Estas comprobaciones corresponden a la implementación actual,
no a una ejecución histórica en 2021. No sustituyen la revisión de los
microdatos originales y de sus etiquetas.

\begin{table}[htbp]
\centering\small
\caption{Estado de la evidencia y de la replicación}\label{tab:estado}
\begin{tabularx}{\textwidth}{@{}YYp{4.0cm}@{}}\toprule
Producto & Qué se ha realizado & Qué falta\\\midrule
Cuestionarios y códigos & Cotejo documental y reglas explícitas & Contrastar distribución en SAV del usuario\\
Gráficos de este artículo & Transcripción verificada del informe & Ningún IC propio calculado\\
Programa R & Pruebas con SPSS artificial y entorno fijado & Ejecutar microdatos autorizados\\
Prevalencias propias & Diseño y exportaciones preparados & Estimar puntos, EE, IC y faltantes\\
Réplica oficial & Referencias CP.5, CP.7 y CP.10 identificadas & Armonizar no respuesta y edades derivadas\\\bottomrule
\end{tabularx}
\fuente{Registro de trabajo del repositorio. ``Pendiente'' no implica un resultado nulo ni un fallo de significación.}
\end{table}

\FloatBarrier
\section{Discusión}
\subsection{Qué permite concluir la evidencia disponible}
Los tabulados oficiales documentan presencia de prácticas violentas y
justificación declarada en las poblaciones cubiertas, así como experiencias
previas de unión temprana. Ofrecen puntos de referencia para la política
pública y para comprobar órdenes de magnitud de una futura réplica. No son
evidencia suficiente para atribuir efectos a pobreza, residencia, educación
o intervenciones.

La lectura conjunta muestra heterogeneidad conceptual. La disciplina violenta
publicada es mayor en el área urbana, mientras la justificación de los cinco
escenarios es mayor en el área rural. Esta comparación entre módulos ayuda
a cuestionar un relato único sobre violencia, pero no estima correlación
individual. Las diferencias de informante, período, composición y declaración
deben mantenerse presentes.

La amplitud de los porcentajes no puede compararse mecánicamente entre
constructos. Disciplina se refiere al último mes y a una combinación de actos;
unión temprana, a una transición acumulada antes de una edad; y actitudes,
a una posición declarada. Un valor numéricamente mayor no implica por sí
solo mayor gravedad ni mayor prioridad relativa sin un marco normativo.

\subsection{Sesgos de medición y límites de cobertura}
El reporte por un adulto puede subestimar prácticas socialmente censuradas o
desconocidas por quien responde. La ventana de un mes omite experiencias
anteriores. Los ítems no registran de forma exhaustiva frecuencia, intensidad,
consecuencias ni contextos de coerción. Una clasificación cero de disciplina
violenta no equivale a una certificación integral de seguridad.

Las respuestas sobre violencia contra la esposa pueden depender del lenguaje
del escenario y del contexto de entrevista. ``No sabe'' puede reflejar
incertidumbre, ambivalencia o dificultad para responder; tratarlo como
rechazo impondría una interpretación no demostrada. La estrategia de casos
clasificables debe acompañarse de faltantes y sensibilidad, no presentarse
como solución automática al sesgo.

La edad declarada de primera unión puede contener errores de recuerdo y
concentración en edades enteras. Restringir a 20 a 24 años evita exposición
incompleta a los umbrales, pero no sesgos de supervivencia y migración.
El hogar observado puede diferir del de origen; riqueza actual no es una
medida inequívoca de condiciones previas a la unión.

\subsection{Comparabilidad y agenda de investigación}
Una réplica debe cotejar \var{MA9} con la edad derivada \var{wagem}, armonizar
las reglas de no respuesta y contrastar universos. La coincidencia aproximada
con un porcentaje redondeado no valida por sí sola todo el procesamiento.
Se requiere revisar pesos, llaves, exclusiones y varianzas, además de resolver
las inconsistencias editoriales detectadas en el informe.

Para estudiar asociaciones entre normas maternas y disciplina se debe
identificar explícitamente madre o cuidadora, establecer un universo vinculado
y considerar selección y temporalidad. Para evaluar políticas se necesitarían
comparaciones temporales o diseños adicionales. El presente protocolo ofrece
una base para formular esas preguntas sin adelantar su respuesta.

Desde el enfoque de derechos, los resultados secundarios justifican examinar
crianza, servicios de apoyo y protección de adolescentes. El marco INSPIRE
orienta ámbitos de acción complementarios, pero los datos aquí presentados
no seleccionan por sí solos una intervención ni cuantifican su efectividad.

\subsection{Una agenda compartida de prevención familiar y ciudadana}
La combinación de fuentes propone un criterio de seguimiento: evaluar
avances por varias dimensiones, sin usar la caída de una tasa como evidencia
de protección universal. La disciplina observada en MICS, las actitudes de
género, la violencia letal y la victimización aportan señales diferentes.
Su articulación permite formular preguntas de política más precisas que una
clasificación global de seguridad o inseguridad.

La recomendación de Infosegura de dar seguimiento a avisos de violencia
intrafamiliar puede dialogar con intervenciones de crianza y prevención de
violencia de género. Se requiere información adicional para establecer
riesgos individuales y rutas de atención: la prevalencia nacional de MICS no
permite identificar personas que realizarán un aviso ni predecir un feminicidio.
Del mismo modo, una pregunta de justificación puede orientar diagnóstico de
normas, pero no debe emplearse para clasificar a una mujer como víctima o agresora.

\begin{table}[htbp]
\centering\small
\caption{Preguntas de política que surgen de la articulación}\label{tab:agenda}
\begin{tabularx}{\textwidth}{@{}YYY@{}}\toprule
Dimensión de MICS & Aporte de Infosegura & Pregunta y evidencia adicional requerida\\\midrule
Disciplina infantil & Enfoque de prevención y oportunidades de desarrollo & ¿Qué apoyos a cuidadores reducen prácticas violentas? Requiere evaluación y seguimiento propios.\\
Justificación de violencia & Seguimiento de avisos y muertes de mujeres & ¿Cómo se relacionan normas, acceso a servicios y respuesta institucional? Requiere datos vinculables y protección de identidad.\\
Unión temprana & Privaciones materiales y uso de espacios & ¿Qué restricciones de autonomía enfrentan las adolescentes? Requiere historia de unión, hogar de origen y trayectorias educativas.\\\bottomrule
\end{tabularx}
\fuente{Síntesis analítica del artículo y recomendaciones de Infosegura, diapositiva 13. Las preguntas no constituyen efectos estimados ni un sistema de predicción individual.}
\end{table}

Una futura investigación podría armonizar rondas de encuesta y territorios
para contrastar estas hipótesis. Tendría que documentar marcos muestrales,
cuestionarios, pesos y períodos. Con la evidencia presente, la recomendación
es conservar trazabilidad y complementar indicadores: no se identifica el
efecto del control del delito, de una reforma o de condiciones económicas
sobre las experiencias de niños y mujeres.


\section{Conclusiones}
La medición reproducible de la violencia familiar exige correspondencia entre
pregunta, código, unidad, período, denominador y ponderador. En la ENS MICS
2014, omitir la selección infantil o confundir convivencia previa con nunca
unión altera el significado del indicador. También es indispensable separar
actitudes de experiencias y reconocer la coexistencia de prácticas.

El artículo presenta evidencia secundaria de tres dimensiones y un protocolo
para estimaciones propias futuras. Los gráficos y tabulados se identifican
como transcripciones del informe oficial; el programa todavía no ha producido
resultados con microdatos reales. Esta distinción permite aprovechar la
información publicada sin simular una réplica empírica inexistente.

El aporte de Infosegura modifica la lectura sustantiva: la reducción de
homicidios generales puede coexistir con señales distintas por género,
exposición a delincuencia y vulnerabilidades económicas. Los indicadores de
MICS aportan una dimensión familiar que ese descenso no permite dar por
resuelta. La articulación respalda una agenda de prevención con seguimiento
de varias fuentes, sin demostrar un tránsito causal entre ellas.

La siguiente etapa consiste en obtener los archivos autorizados, revisar su
diccionario, ejecutar el diseño complejo y documentar concordancias y
diferencias con la referencia oficial. Hasta entonces, las conclusiones
causales y los contrastes de significación permanecen fuera del alcance del
trabajo.

\section*{Disponibilidad y resguardo de los datos}
El código y la documentación se encuentran en
\url{https://github.com/vtablas001/Violence_MICS}. No se redistribuyen microdatos
ni registros individuales. Las condiciones del proveedor rigen su acceso y
uso; las tablas que se generen deberán someterse a revisión de divulgación.
Las fuentes se consultaron en 2026 y se citan por su fecha original. La fecha
editorial de referencia no altera las fechas de ejecución o de publicación
de los archivos del repositorio.

\begin{thebibliography}{9}
\bibitem[Lumley(2004)]{lumley}
Lumley, T. (2004).
Analysis of Complex Survey Samples.
\emph{Journal of Statistical Software}, 9(8), 1-19.
\url{https://doi.org/10.18637/jss.v009.i08}.

\bibitem[PNUD/Infosegura(2021)]{infosegura2021}
PNUD a través de Infosegura (2021).
\emph{Análisis Multidimensional de Seguridad Ciudadana. Primer Semestre 2021.
El Salvador}. Presentación de DATACCIÓN, versión 3.0 facilitada por el usuario.
Diapositivas 4, 5, 7, 8, 10, 12 y 13. Se utiliza esta versión concreta y no
se le atribuye una URL pública no verificada. Los pies de las diapositivas
identifican a los productores originales de cada serie.

\bibitem[MINSAL-INS, DIGESTYC y UNICEF(s.\,f.)]{ens}
Ministerio de Salud, Instituto Nacional de Salud, Dirección General de
Estadísticas y Censos y UNICEF (s.\,f.).
\emph{Encuesta Nacional de Salud de Indicadores Múltiples por Conglomerados
2014: informe final}. Tablas CP.5, CP.7 y CP.10; nota 5, p. 306; apéndices A y F.
\url{https://www.unicef.org/elsalvador/media/1261/file/mics_final.pdf}.
Se identifica el año de encuesta sin atribuir una fecha editorial no verificada.

\bibitem[MINSAL-INS, DIGESTYC y UNICEF(2014)]{catalogo}
Ministerio de Salud, Instituto Nacional de Salud, Dirección General de
Estadísticas y Censos y UNICEF (2014).
\emph{Encuesta de Indicadores Múltiples por Conglomerados, El Salvador}.
Referencia \var{SLV_2014_MICS_v01_M}. Metadatos catalogados en 2017.
\url{https://microdata.worldbank.org/index.php/catalog/2780}.

\bibitem[ONU Mujeres, UNICEF y UNFPA(2017)]{onu2017}
ONU Mujeres, UNICEF y UNFPA (2017, 17 de agosto).
\emph{Agencias de las Naciones Unidas reconocen importantes reformas
legislativas para prohibir el matrimonio infantil y las uniones tempranas en
El Salvador}.
\href{https://www.unicef.org/elsalvador/comunicados-prensa/agencias-de-las-naciones-unidas-reconocen-importantes-reformas-legislativas-para}{Comunicado oficial de UNICEF El Salvador}.

\bibitem[OMS y asociados(2016)]{inspire}
Organización Mundial de la Salud y asociados (2016).
\emph{INSPIRE: Seven strategies for ending violence against children}.
Ginebra: OMS.
\href{https://www.who.int/publications/i/item/inspire-seven-strategies-for-ending-violence-against-children}{Publicación oficial}.
\end{thebibliography}

\clearpage
\appendix
\section{Diccionario mínimo y reglas de validación}\label{anexo:diccionario}
Los nombres se limpian a minúsculas al leer los archivos. El cuadro conserva
su forma original para facilitar el cotejo con el cuestionario. Los códigos
de no sabe o falta de respuesta no se interpretan como respuestas sustantivas.

\begin{longtable}{@{}p{2.4cm}p{2.6cm}p{9.9cm}@{}}
\caption{Campos requeridos y comprobaciones}\label{tab:diccionario}\\
\toprule Archivo & Campo & Definición y regla\\\midrule\endfirsthead
\toprule Archivo & Campo & Definición y regla (continuación)\\\midrule\endhead
\bottomrule\endfoot
hh, hl, wm & HH1, HH2 & Identificadores del hogar. No se admiten llaves faltantes ni hogares duplicados.\\
hl & HL1 & Línea de miembro; única dentro de hogar.\\
wm & WM4 & Línea de mujer; única dentro de hogar.\\
hh & HH9 & Resultado de entrevista; se requieren hogares completos, código 1.\\
wm & WM7 & Resultado de entrevista; se requieren mujeres completas, código 1.\\
hh & SL1 & Conteo entero no negativo de miembros 1-17; debe coincidir con HL6.\\
hh & SL9B, SL9C & Línea y edad del seleccionado; deben coincidir con el listado.\\
hh & CD2 & Línea del niño 1-14; debe ser la misma que SL9B.\\
hl & HL4, HL6 & Sexo y edad cumplida. El sexo se presenta como factor.\\
hh & HH6 & Área: 1 urbana, 2 rural.\\
hh & windex5 & Quintil nacional combinado, 1 a 5.\\
hh & psu, stratum & UPM y estrato; no se sustituyen por región.\\
hh & hhweight & Ponderador positivo; se multiplica por SL1 para infancia.\\
wm & wmweight & Ponderador individual positivo, sin sustitución.\\
hh & CD3A-K & Sí 1, no 2; otros códigos documentados como desconocidos se conservan como NA.\\
wm & DV1A-J & Sí 1, no 2, no sabe 8; faltantes permanecen desconocidos.\\
wm & WB2 & Edad actual; entrevistas de mujeres 15-49.\\
wm & MA1 & 1 casada; 2 acompañada; 3 no en unión actual.\\
wm & MA5 & Si MA1=3: 1 casada antes; 2 acompañada antes; 3 nunca unida.\\
wm & MA9 & Edad declarada de primera unión; 98/99 desconocidos en el protocolo. Se rechazan edades futuras o incompatibles con trayectoria.\\
\end{longtable}

\section{Trazabilidad de las cifras secundarias}\label{anexo:trazabilidad}
Los datos de los gráficos se guardan como archivos CSV dentro del proyecto.
Son agregados transcritos del informe, no microdatos. Las páginas de la columna
final corresponden a la numeración impresa; la página del visor PDF es dos
unidades mayor. El compuesto de diez escenarios no se calcula con marginales.

\begin{table}[htbp]
\centering\small
\caption{Origen verificable de los valores graficados}\label{tab:traza}
\begin{tabularx}{\textwidth}{@{}p{2.7cm}Y>{\raggedright\arraybackslash}p{2.8cm}@{}}\toprule
Archivo o elemento & Valores porcentuales y dominio & Localizador\\\midrule
\var{disciplina.csv} & 35,5; 32,2; 39,4; 3,0; 52,0. Niños 1-14, total nacional & CP.5, p. 282; PDF 284\\
\var{actitudes.csv} & A-J: 1,6; 5,2; 2,5; 1,8; 2,0; 5,6; 2,4; 0,7; 1,3; 1,3. Mujeres 15-49 & CP.10, p. 296; PDF 298\\
\var{union.csv} & 5,8 y 25,5; primera unión antes de 15 y 18, mujeres 20-24 & CP.7, p. 288; PDF 290\\
Cuadro \ref{tab:discgrupo} & Filas Total, Masculino, Femenino, Urbana y Rural & CP.5, p. 282\\
Cuadro \ref{tab:actgrupos} & Compuesto A-E: Total, Urbana, Rural y 20-24 & CP.10, p. 296\\
Cuadro \ref{tab:union} & Totales 15-49 y 20-49; fila 20-24 & CP.7, pp. 287-288\\\bottomrule
\end{tabularx}
\end{table}

\begin{samepage}
La fuente descargada se identifica mediante SHA256:
\begin{quote}\begin{minipage}{\linewidth}\footnotesize\ttfamily
81eac967646e1808434c5dd21c9d30237\\
eaea93c33795ea6b281e10709e4c1a0
\end{minipage}\end{quote}
\end{samepage}
El salto de línea divide visualmente una sola cadena hexadecimal. No se
reconstruyen tamaños sin ponderar a partir de las bases ponderadas publicadas
ni se fabrican intervalos. Las diferencias mencionadas en el texto son restas
de porcentajes redondeados y se identifican como tales.

\section{Sensibilidad a desconocidos: ejemplo analítico}
Si la fracción ponderada desconocida es $m$ y la proporción entre casos
clasificables es $p_O$, los límites extremos son
$L=(1-m)p_O$ y $U=(1-m)p_O+m$. La figura \ref{fig:sensibilidad} fija
$p_O=0{,}40$ únicamente para ilustrar la identidad. Ninguno de sus puntos es
una estimación de El Salvador ni una simulación calibrada al país.

\begin{figure}[htbp]
\centering
\begin{tikzpicture}
\begin{axis}[width=0.86\textwidth,height=6cm,xmin=0,xmax=30,ymin=20,ymax=65,
 xlabel={Fracción ponderada desconocida (\%)},ylabel={Proporción posible (\%)},
 grid=major,grid style={gray!15},legend style={font=\small,at={(0.5,-0.25)},anchor=north,legend columns=2}]
\addplot[petroleo,thick,domain=0:30,samples=31]{40*(1-x/100)};
\addlegendentry{Límite inferior}
\addplot[ocre,thick,domain=0:30,samples=31]{40*(1-x/100)+x};
\addlegendentry{Límite superior}
\end{axis}
\end{tikzpicture}
\caption{Amplitud por información faltante: identidad bajo un ejemplo hipotético}\label{fig:sensibilidad}
\fuente{Cálculo algebraico ilustrativo con $p_O=0{,}40$. No son datos empíricos, IC ni resultados del programa aplicado a la encuesta.}
\end{figure}

El ejemplo muestra que ignorar faltantes puede ocultar un rango amplio de
valores compatibles con la información observada. No prueba que los
desconocidos tengan mayor o menor prevalencia, ni reemplaza estrategias de
imputación sustentadas en supuestos explícitos.

\FloatBarrier
\section{Protocolo de ejecución y tabulado propio pendiente}
La biblioteca de R se fija en \var{renv.lock}. La secuencia reproducible es:
obtener los tres SAV bajo condiciones autorizadas; restaurar dependencias;
ejecutar pruebas; correr el análisis en una carpeta vacía; revisar el
diccionario observado, la auditoría de diseño y los estados de estimación;
y cotejar definiciones antes de comparar con los tabulados del informe.

\begin{verbatim}
Rscript -e 'renv::load(); source("tests/run_tests.R")'
Rscript -e 'renv::load(); source("scripts/run_analysis.R")' \
  data/raw results/real
\end{verbatim}
El salto con barra invertida corresponde a una terminal compatible; en
PowerShell se puede escribir el segundo comando en una sola línea. En R,
la primera preparación consiste en instalar \var{renv}, cargar el proyecto
y ejecutar \var{renv::restore(prompt=FALSE)}.

\begin{table}[htbp]
\centering\small
\caption{Estructura del tabulado propio a completar con microdatos}\label{tab:plantilla}
\begin{tabularx}{\textwidth}{@{}Yccccc@{}}\toprule
Indicador y dominio & $n$ válido & \% & EE & IC95 & \% desc.\\\midrule
Disciplina violenta, 1-14 & --- & --- & --- & --- & ---\\
Físico severo, 1-14 & --- & --- & --- & --- & ---\\
Justifica cinco, 15-49 & --- & --- & --- & --- & ---\\
Justifica diez, 15-49 & --- & --- & --- & --- & ---\\
Unión antes de 15, 20-24 & --- & --- & --- & --- & ---\\
Unión antes de 18, 20-24 & --- & --- & --- & --- & ---\\\bottomrule
\end{tabularx}
\fuente{Plantilla de resultados propios pendientes. Los guiones significan ``no calculado'', nunca cero. No copiar a estas celdas cifras oficiales como si fueran reestimaciones.}
\end{table}

El archivo de estimaciones incorpora además límites extremos por desconocidos,
grados de libertad, número de afirmativas y estado de la estimación. La
salida de auditoría identifica filas del diseño, UPM, estratos, estratos con
una UPM y rango de pesos. Las huellas de archivos y la información de sesión
permiten identificar la versión local usada.

\section{Correcciones y criterios de lectura}
\begin{table}[htbp]
\centering\small
\caption{Correcciones que cambian el significado del análisis}\label{tab:correcciones}
\begin{tabularx}{\textwidth}{@{}YY@{}}\toprule
Riesgo de interpretación & Regla adoptada\\\midrule
Usar MA5=2 como nunca unida & Nunca unida requiere MA1=3 y MA5=3\\
Ponderar infancia solo con hhweight & Multiplicar por SL1, miembros 1-17\\
Exigir nombres con letras mayúsculas después de limpiar & Usar sufijos minúsculos coherentes\\
Filtrar edades antes de crear diseño & Construir diseño y luego estimar dominios\\
Tratar algún método no violento como exclusivo & Exigir además ausencia de todo acto violento\\
Sumar los diez porcentajes DV & Construir unión de eventos en microdatos\\
Presentar no sabe como no & Conservar desconocidos y sus denominadores\\
Leer diferencias de porcentajes como efectos & Mantener interpretación descriptiva e incertidumbre\\\bottomrule
\end{tabularx}
\end{table}

Los gráficos oficiales, la ilustración matemática y las estimaciones propias
pendientes mantienen identificaciones separadas para facilitar su lectura
y replicación.

\section{Procedencia y reproducción de las series de seguridad}\label{anexo:infosegura}
Las seis series de PNUD/Infosegura (2021) complementan MICS con indicadores
de seguridad y condiciones de vida. Los CSV, el manifiesto de procedencia y
el programa R permiten repetir los cálculos sobre agregados secundarios,
sin registros individuales.

\begin{table}[htbp]
\centering\small
\caption{Series de seguridad: archivos, dimensiones y restricciones}\label{tab:traza_infosegura}
\begin{tabularx}{\textwidth}{@{}p{3.5cm}p{2.5cm}Y@{}}\toprule
Archivo, prefijo \var{infosegura_} & Dimensión & Alcance y restricción\\\midrule
\var{homicidios.csv} & Violencia letal & Tasas por población de referencia; 2021 es proyección.\\
\var{victimizacion.csv} & Experiencia delictiva & Delincuencia en doce meses; universo etario no explicitado.\\
\var{preocupaciones.csv} & Opinión & COVID-19 solo consta en 2021; antes permanece vacío, no cero.\\
\var{denuncias.csv} & Demanda institucional & Promedios mensuales; corte de 2021 no precisado.\\
\var{privaciones.csv} & Condiciones de vida & Cambios 2019--2020 en puntos porcentuales, no niveles.\\
\var{percepcion.csv} & Seguridad percibida & Transporte y comunidad; distinto de uso de espacios.\\\bottomrule
\end{tabularx}
\end{table}

El manifiesto \var{data/infosegura_fuente.json} registra procedencia,
localizadores y versión. Los CSV conservan la precisión disponible y las
figuras no imputan años ausentes. \var{rebuild_analysis.R} regenera los contrastes.

\begin{samepage}
Reproducción desde la carpeta del proyecto, con R base:
\begin{verbatim}
Rscript --vanilla rebuild_analysis.R
\end{verbatim}
\end{samepage}
Se mantienen las series disponibles en 2021 y se separan proyecciones de
observaciones. Migración, revictimización y series mixtas se excluyen por
bases insuficientemente verificadas. No se calculan IC propios sobre estos
agregados ni se certifica comparabilidad de ponderadores entre rondas.


\end{document}

